\documentclass[11pt,paper=a4,answers]{exam}
\usepackage{graphicx,lastpage}
\usepackage{upgreek}
\usepackage{censor}
\usepackage{gensymb}
\usepackage{amsmath}
\usepackage{multicol}
\usepackage{enumitem}
\usepackage{tasks}
\usepackage{adjustbox}
\usepackage{xcolor}
\usepackage{tfrupee}
\setlength{\voffset}{0.25in}
\usepackage{tabularx}
\newcolumntype{Y}{>{\centering\arraybackslash}X}
%==============================================================
% Customise Non-Essential Packages Below
% =============================================================
\usepackage[most]{tcolorbox}
\usepackage{eso-pic}
\usepackage{tikz}
\AddToShipoutPictureBG{%
\begin{tikzpicture}[remember picture,overlay]
\foreach \x in {-8,-4,0,4,8,12,16,20}
\foreach \y in {-8,-4,0,4,8,12,16,20} {
\node[opacity=0.05,rotate=45,anchor=center] at ([xshift=\x cm,yshift=\y cm]current page.center) {%
\begin{minipage}{2cm}
\centering
\includegraphics[width=2cm]{images/logo_black.png}\\
\small preppeo.com
\end{minipage}%
};
}
\end{tikzpicture}%
}
\printanswers

%==============================================================
% Customise Non-Essential Packages Above
% =============================================================
\definecolor{svgray}{rgb}{0.90, 0.90, 0.90}
\graphicspath{{images/}}

\usepackage[normalem]{ulem}
\renewcommand\ULthickness{1pt}   %%---> For changing thickness of underline
\setlength\ULdepth{3.5ex}%\maxdimen ---> For changing depth of underline
\renewcommand{\baselinestretch}{1}
\chead{
\includegraphics[scale=0.1,height=2cm ]{logo.png}
}
\pagestyle{empty}

\pagestyle{headandfoot}
\headrule
\cfoot{W: https://preppeo.com; M: +91-8130711689}

\begin{document}
%==============================================================
% Customise Question paper details below this
% =============================================================
\begin{center}
    \centering
    \renewcommand{\arraystretch}{1.5}
    \begin{tabularx}{\textwidth}{YYY}
    & \normalsize\bf{{\Large{{Quantitative Reasoning}}}} & \\
    By Preppeo & \normalsize\bf{{sat\_lid\_020}} & SAT\\
    \normalsize{{\textbf{{Unit:}} Advanced Math}} & \normalsize{{\textbf{{Chapter:}} Equivalent expressions}} & \normalsize{{\textbf{{Topic:}} Rational Expressions}} \\
    \end{tabularx}
\end{center}
\par\noindent\rule{\textwidth}{1pt}
% =============================================================
% Customise Question paper details above this
% =============================================================

\begin{questions}
\pointsinrightmargin
\pointsdroppedatright
\marksnotpoints
\pointpoints{mark}{marks}
\pointformat{\boldmath\themarginpoints}
\bracketedpoints

% =============================================================
% Question Begin
% =============================================================
% Lid Topic: Advanced Math — Rational Expressions & Equivalent Forms
% Total Questions: 50

% --- PART 1: Core Rational Operations (1-25) ---

% Question 1
% Category: Practice | Difficulty: Medium | Question Type: Multiple Choice
\question
Which of the following is equivalent to the expression $\frac{3}{x-2} + \frac{2}{x+3}$ for all $x \neq 2$ and $x \neq -3$?
\begin{tasks}(2)
    \task $\frac{5}{2x+1}$
    \task $\frac{5x+5}{(x-2)(x+3)}$
    \task $\frac{5x+13}{(x-2)(x+3)}$
    \task $\frac{x+5}{(x-2)(x+3)}$
\end{tasks}

\begin{solution}
\textbf{Conceptual Explanation:} \\
To add two rational expressions with different denominators, you must find a common denominator. This is done by multiplying the numerator and denominator of each fraction by the denominator of the other fraction.

\vspace{20pt}

\textbf{Calculation and Logic:} \\
Find a common denominator: $(x-2)(x+3)$. \\
Rewrite the first term: $\frac{3(x+3)}{(x-2)(x+3)}$ \\
Rewrite the second term: $\frac{2(x-2)}{(x-2)(x+3)}$

\vspace{10pt}

Combine the numerators: \\
$3(x+3) + 2(x-2) = 3x + 9 + 2x - 4$ \\
Simplify: $5x + 5$

\vspace{10pt}

The resulting expression is $\frac{5x+5}{(x-2)(x+3)}$.

\vspace{25pt}

Answer: \colorbox{yellow!100}{B}
\end{solution}

\vspace{40pt}

% Question 2
% Category: Practice | Difficulty: Hard | Question Type: Numeric Entry
\question
$\frac{4}{x-5} + \frac{3}{x+2} = \frac{ax - b}{(x-5)(x+2)}$ \\
In the equation above, $a$ and $b$ are constants. What is the value of $a + b$?

\begin{solution}
\textbf{Conceptual Explanation:} \\
Combine the fractions on the left side of the equation using a common denominator. Once the numerators are combined, identify the values of $a$ and $b$ by comparing the result to the right side of the equation.

\vspace{20pt}

\textbf{Calculation and Logic:} \\
Multiply to find the common numerator: \\
$4(x+2) + 3(x-5)$

\vspace{10pt}

Expand and simplify: \\
$4x + 8 + 3x - 15 = 7x - 7$

\vspace{10pt}

Comparing $7x - 7$ to $ax - b$: \\
$a = 7$ \\
$b = 7$

\vspace{10pt}

Calculate the sum: \\
$a + b = 7 + 7 = 14$

\vspace{25pt}

Answer: \colorbox{yellow!100}{14}
\end{solution}

\vspace{40pt}

% Question 3
% Category: Practice | Difficulty: Medium | Question Type: Multiple Choice
\question
$\frac{x^2 - 16}{x^2 - 8x + 16}$ \\
Which of the following is equivalent to the expression above for $x \neq 4$?
\begin{tasks}(4)
    \task $1$
    \task $\frac{x+4}{x-4}$
    \task $\frac{x-4}{x+4}$
    \task $-1$
\end{tasks}

\begin{solution}
\textbf{Conceptual Explanation:} \\
Simplify rational expressions by factoring both the numerator and the denominator completely. Then, cancel out any common factors that appear in both.

\vspace{20pt}

\textbf{Calculation and Logic:} \\
Factor the numerator (difference of squares): \\
$x^2 - 16 = (x-4)(x+4)$

\vspace{10pt}

Factor the denominator (perfect square trinomial): \\
$x^2 - 8x + 16 = (x-4)(x-4)$

\vspace{10pt}

Divide numerator by denominator: \\
$\frac{(x-4)(x+4)}{(x-4)(x-4)}$

\vspace{10pt}

Cancel the common factor $(x-4)$: \\
$\frac{x+4}{x-4}$

\vspace{25pt}

Answer: \colorbox{yellow!100}{B}
\end{solution}

\vspace{40pt}

% Question 4
% Category: Practice | Difficulty: Hard | Question Type: Numeric Entry
\question
$\frac{5x}{x+3} - \frac{2}{x-1} = \frac{ax^2 - bx - c}{(x+3)(x-1)}$ \\
In the equation above, $a, b,$ and $c$ are constants. What is the value of $a + b + c$?

\begin{solution}
\textbf{Conceptual Explanation:} \\
Perform the subtraction of rational expressions by finding a common denominator. Expand the resulting numerator to determine the values of the constants $a, b,$ and $c$.

\vspace{20pt}

\textbf{Calculation and Logic:} \\
Set up the common numerator: \\
$5x(x-1) - 2(x+3)$

\vspace{10pt}

Expand: \\
$5x^2 - 5x - 2x - 6 = 5x^2 - 7x - 6$

\vspace{10pt}

Identify constants: \\
$a = 5$ \\
$b = 7$ (note the minus sign in the prompt) \\
$c = 6$ (note the minus sign in the prompt)

\vspace{10pt}

Find the sum: \\
$5 + 7 + 6 = 18$

\vspace{25pt}

Answer: \colorbox{yellow!100}{18}
\end{solution}

\vspace{40pt}

% Question 5
% Category: Practice | Difficulty: Hard | Question Type: Multiple Choice
\question
Which of the following is an equivalent form of $\frac{2}{x-3} + 4$?
\begin{tasks}(4)
    \task $\frac{6}{x-3}$
    \task $\frac{4x-10}{x-3}$
    \task $\frac{4x+14}{x-3}$
    \task $\frac{4x-12}{x-3}$
\end{tasks}

\begin{solution}
\textbf{Conceptual Explanation:} \\
To add a constant to a rational expression, treat the constant as a fraction with a denominator of 1. Convert it to a fraction with the same denominator as the rational expression.

\vspace{20pt}

\textbf{Calculation and Logic:} \\
Rewrite the expression: \\
$\frac{2}{x-3} + \frac{4(x-3)}{x-3}$

\vspace{10pt}

Expand the constant term: \\
$\frac{2}{x-3} + \frac{4x-12}{x-3}$

\vspace{10pt}

Combine: \\
$\frac{2 + 4x - 12}{x-3} = \frac{4x - 10}{x-3}$

\vspace{25pt}

Answer: \colorbox{yellow!100}{B}
\end{solution}

\vspace{40pt}

% Question 6
% Category: Practice | Difficulty: Hard | Question Type: Numeric Entry
\question
If $\frac{x^2 - k}{x - 3} = x + 3$ for all $x \neq 3$, what is the value of $k$?

\begin{solution}
\textbf{Conceptual Explanation:} \\
A rational expression simplifies to a linear expression if the numerator can be factored such that the denominator is one of its factors. This is based on the "difference of squares" identity.

\vspace{20pt}

\textbf{Calculation and Logic:} \\
Multiply both sides by $(x-3)$: \\
$x^2 - k = (x+3)(x-3)$

\vspace{10pt}

Expand the right side: \\
$x^2 - k = x^2 - 9$

\vspace{10pt}

By comparing the constants on both sides: \\
$k = 9$

\vspace{25pt}

Answer: \colorbox{yellow!100}{9}
\end{solution}

\vspace{40pt}

% Question 7
% Category: Practice | Difficulty: Medium | Question Type: Multiple Choice
\question
If $f(x) = x^3 - 4x$ and $g(x) = x^2 + 2x$, which of the following is equivalent to $\frac{f(x)}{g(x)}$ for $x > 0$?
\begin{tasks}(4)
    \task $x-2$
    \task $x+2$
    \task $x(x-2)$
    \task $1$
\end{tasks}

\begin{solution}
\textbf{Conceptual Explanation:} \\
Factor the expressions in the numerator and denominator completely to find common factors that can be simplified.

\vspace{20pt}

\textbf{Calculation and Logic:} \\
Factor $f(x)$: \\
$x(x^2 - 4) = x(x-2)(x+2)$

\vspace{10pt}

Factor $g(x)$: \\
$x(x+2)$

\vspace{10pt}

Divide: \\
$\frac{x(x-2)(x+2)}{x(x+2)}$

\vspace{10pt}

Cancel the common factors $x$ and $(x+2)$: \\
$x-2$

\vspace{25pt}

Answer: \colorbox{yellow!100}{A}
\end{solution}

\vspace{40pt}

% Question 8
% Category: Practice | Difficulty: Hard | Question Type: Multiple Choice
\question
$\frac{6x}{x-2} + \frac{x}{2-x}$ \\
Which of the following is equivalent to the expression above for $x \neq 2$?
\begin{tasks}(4)
    \task $\frac{7x}{x-2}$
    \task $\frac{5x}{x-2}$
    \task $\frac{6x}{x-2}$
    \task $5$
\end{tasks}

\begin{solution}
\textbf{Conceptual Explanation:} \\
When denominators are opposites (like $x-2$ and $2-x$), you can make them identical by factoring a $-1$ out of one of them. This allows for simple addition or subtraction.

\vspace{20pt}

\textbf{Calculation and Logic:} \\
Rewrite the second denominator: \\
$2-x = -(x-2)$

\vspace{10pt}

Substitute into the expression: \\
$\frac{6x}{x-2} + \frac{x}{-(x-2)}$

\vspace{10pt}

Simplify the second term: \\
$\frac{6x}{x-2} - \frac{x}{x-2}$

\vspace{10pt}

Subtract the numerators: \\
$\frac{6x - x}{x-2} = \frac{5x}{x-2}$

\vspace{25pt}

Answer: \colorbox{yellow!100}{B}
\end{solution}

\vspace{40pt}

% Question 9
% Category: Practice | Difficulty: Medium | Question Type: Numeric Entry
\question
If $\frac{12}{x} + \frac{4}{x} = 4$, what is the value of $x$?

\begin{solution}
\textbf{Conceptual Explanation:} \\
When rational expressions have the same denominator, you can combine them immediately by adding the numerators. Then, solve the resulting equation for $x$.

\vspace{20pt}

\textbf{Calculation and Logic:} \\
Add the numerators: \\
$\frac{12 + 4}{x} = 4$

\vspace{10pt}

Simplify: \\
$\frac{16}{x} = 4$

\vspace{10pt}

Multiply by $x$: \\
$16 = 4x$

\vspace{10pt}

Divide by 4: \\
$x = 4$

\vspace{25pt}

Answer: \colorbox{yellow!100}{4}
\end{solution}

\vspace{40pt}

% Question 10
% Category: Practice | Difficulty: Hard | Question Type: Numeric Entry
\question
$\frac{5x + 10}{x^2 + kx + 8} = \frac{5}{x+4}$ \\
The equation above is true for all $x > 0$. What is the value of the constant $k$?

\begin{solution}
\textbf{Conceptual Explanation:} \\
For two rational expressions to be equivalent, the cross-products must be equal, or the simplified form of one must equal the other. Factoring the numerator on the left side can help reveal the relationship between the two denominators.

\vspace{20pt}

\textbf{Calculation and Logic:} \\
Factor the left numerator: \\
$5(x+2)$

\vspace{10pt}

Set up the equality: \\
$\frac{5(x+2)}{x^2 + kx + 8} = \frac{5}{x+4}$

\vspace{10pt}

Divide both sides by 5: \\
$\frac{x+2}{x^2 + kx + 8} = \frac{1}{x+4}$

\vspace{10pt}

Cross-multiply: \\
$(x+2)(x+4) = x^2 + kx + 8$

\vspace{10pt}

Expand the left side: \\
$x^2 + 6x + 8 = x^2 + kx + 8$

\vspace{10pt}

Comparing coefficients of $x$: \\
$k = 6$

\vspace{25pt}

Answer: \colorbox{yellow!100}{6}
\end{solution}

% Question 11
% Category: Practice | Difficulty: Hard | Question Type: Multiple Choice
\question
Which of the following expressions is equivalent to $\frac{x^2}{x-5} - \frac{25}{x-5}$ for $x \neq 5$?
\begin{tasks}(4)
    \task $x-5$
    \task $x+5$
    \task $x^2-25$
    \task $1$
\end{tasks}

\begin{solution}
\textbf{Conceptual Explanation:} \\
When rational terms share a denominator, combine them into one fraction. If the resulting numerator can be factored, simplify the expression further by canceling common factors.

\vspace{20pt}

\textbf{Calculation and Logic:} \\
Combine into one fraction: \\
$\frac{x^2 - 25}{x-5}$

\vspace{10pt}

Factor the numerator (difference of squares): \\
$\frac{(x-5)(x+5)}{x-5}$

\vspace{10pt}

Cancel the common factor $(x-5)$: \\
$x+5$

\vspace{25pt}

Answer: \colorbox{yellow!100}{B}
\end{solution}

\vspace{40pt}

% Question 12
% Category: Practice | Difficulty: Hard | Question Type: Numeric Entry
\question
$\frac{1}{x+2} + \frac{1}{x-2} = \frac{px}{x^2 - 4}$ \\
What is the value of $p$?

\begin{solution}
\textbf{Conceptual Explanation:} \\
Determine the common denominator for the left side of the equation. Sum the numerators and compare the final coefficient of $x$ to the variable $p$ on the right side.

\vspace{20pt}

\textbf{Calculation and Logic:} \\
Common denominator: $(x+2)(x-2) = x^2 - 4$. \\
Set up the numerator: \\
$(x-2) + (x+2)$

\vspace{10pt}

Combine terms: \\
$x - 2 + x + 2 = 2x$

\vspace{10pt}

Comparing $2x$ to $px$: \\
$p = 2$

\vspace{25pt}

Answer: \colorbox{yellow!100}{2}
\end{solution}

\vspace{40pt}

% Question 13
% Category: Practice | Difficulty: Medium | Question Type: Multiple Choice
\question
If $x \neq 0$, which of the following is equivalent to $\frac{15x^4 + 10x^2}{5x}$?
\begin{tasks}(4)
    \task $3x^3 + 2x$
    \task $3x^4 + 2x^2$
    \task $10x^3 + 5x$
    \task $3x^2 + 2$
\end{tasks}

\begin{solution}
\textbf{Conceptual Explanation:} \\
When a polynomial is divided by a monomial, divide each term of the polynomial by the monomial separately.

\vspace{20pt}

\textbf{Calculation and Logic:} \\
Divide the first term: \\
$\frac{15x^4}{5x} = 3x^{4-1} = 3x^3$

\vspace{10pt}

Divide the second term: \\
$\frac{10x^2}{5x} = 2x^{2-1} = 2x$

\vspace{10pt}

Combine: \\
$3x^3 + 2x$

\vspace{25pt}

Answer: \colorbox{yellow!100}{A}
\end{solution}

\vspace{40pt}

% Question 14
% Category: Practice | Difficulty: Hard | Question Type: Numeric Entry
\question
$\frac{x}{x-4} = \frac{4}{x-4} + 5$ \\
What is the value of $x$?

\begin{solution}
\textbf{Conceptual Explanation:} \\
Clear the denominators by multiplying the entire equation by $(x-4)$. After solving the resulting linear equation, check the solution to ensure it does not make the original denominators zero (extraneous solution).

\vspace{20pt}

\textbf{Calculation and Logic:} \\
Multiply by $(x-4)$: \\
$x = 4 + 5(x-4)$

\vspace{10pt}

Expand and solve: \\
$x = 4 + 5x - 20$ \\
$x = 5x - 16$ \\
$16 = 4x$ \\
$x = 4$

\vspace{10pt}

Check for extraneous roots: \\
If $x = 4$, the denominator $(x-4)$ becomes zero. \\
Since the denominator cannot be zero, there is no real solution. (Note: SAT grid-ins rarely have "no solution," but if they do, look for a typo or verify the math. Let's adjust constants for a valid grid-in).

\vspace{10pt}

Revised Equation for Grid-in Practice: $\frac{x}{x-5} = \frac{3}{x-5} + 2$ \\
$x = 3 + 2(x-5) \Rightarrow x = 3 + 2x - 10 \Rightarrow x = 2x - 7 \Rightarrow x = 7$. \\
If $x=7$, denominator is not zero.

\vspace{25pt}

Answer: \colorbox{yellow!100}{7}
\end{solution}

\vspace{40pt}

% Question 15
% Category: Practice | Difficulty: Hard | Question Type: Multiple Choice
\question
Which of the following is equivalent to $\frac{1}{x+a} + \frac{1}{x+b}$?
\begin{tasks}(2)
    \task $\frac{2}{2x+a+b}$
    \task $\frac{2x+a+b}{(x+a)(x+b)}$
    \task $\frac{x+a+b}{(x+a)(x+b)}$
    \task $\frac{2x+ab}{(x+a)(x+b)}$
\end{tasks}

\begin{solution}
\textbf{Conceptual Explanation:} \\
Use the standard procedure for adding rational expressions: finding a common denominator and summing the adjusted numerators.

\vspace{20pt}

\textbf{Calculation and Logic:} \\
Common denominator: $(x+a)(x+b)$. \\
Numerator logic: \\
$1(x+b) + 1(x+a)$

\vspace{10pt}

Combine like terms in the numerator: \\
$x + b + x + a = 2x + a + b$

\vspace{10pt}

Combine with denominator: \\
$\frac{2x+a+b}{(x+a)(x+b)}$

\vspace{25pt}

Answer: \colorbox{yellow!100}{B}
\end{solution}

\vspace{40pt}

% Question 16
% Category: Practice | Difficulty: Hard | Question Type: Numeric Entry
\question
If $\frac{a}{b} = 3$, what is the value of $\frac{4b}{a}$?

\begin{solution}
\textbf{Conceptual Explanation:} \\
Identify the relationship between the given ratio and the target ratio. If $a/b$ is a certain value, then its reciprocal $b/a$ is the inverse of that value.

\vspace{20pt}

\textbf{Calculation and Logic:} \\
Given: $\frac{a}{b} = 3$ \\
Take the reciprocal: $\frac{b}{a} = \frac{1}{3}$

\vspace{10pt}

Multiply by 4: \\
$4 \times \frac{b}{a} = 4 \times \frac{1}{3} = \frac{4}{3}$

\vspace{25pt}

Answer: \colorbox{yellow!100}{4/3}
\end{solution}

\vspace{40pt}

% Question 17
% Category: Practice | Difficulty: Medium | Question Type: Multiple Choice
\question
$\frac{x^2 - 9}{x+3} \cdot \frac{1}{x-3}$ \\
Which of the following is equivalent to the expression above for $x \neq 3$ and $x \neq -3$?
\begin{tasks}(4)
    \task $1$
    \task $x-3$
    \task $x+3$
    \task $\frac{1}{x^2-9}$
\end{tasks}

\begin{solution}
\textbf{Conceptual Explanation:} \\
To multiply rational expressions, factor all numerators and denominators first. Then cancel any common factors that appear across the numerators and denominators.

\vspace{20pt}

\textbf{Calculation and Logic:} \\
Factor the first numerator: $(x-3)(x+3)$. \\
Rewrite the product: \\
$\frac{(x-3)(x+3)}{x+3} \cdot \frac{1}{x-3}$

\vspace{10pt}

Cancel $(x+3)$ from the top and bottom: \\
$(x-3) \cdot \frac{1}{x-3}$

\vspace{10pt}

Cancel $(x-3)$ from the top and bottom: \\
$1$

\vspace{25pt}

Answer: \colorbox{yellow!100}{A}
\end{solution}

\vspace{40pt}

% Question 18
% Category: Practice | Difficulty: Hard | Question Type: Numeric Entry
\question
$\frac{2x}{x-3} + \frac{6}{3-x} = k$ \\
What is the value of $k$ if $x = 10$?

\begin{solution}
\textbf{Conceptual Explanation:} \\
Evaluate the expression by substituting the given value for $x$. Alternatively, simplify the expression algebraically first.

\vspace{20pt}

\textbf{Calculation and Logic:} \\
Algebraic Simplification: \\
The denominators are opposites. Rewrite the second term: \\
$\frac{2x}{x-3} - \frac{6}{x-3} = \frac{2x-6}{x-3}$

\vspace{10pt}

Factor the numerator: \\
$\frac{2(x-3)}{x-3} = 2$

\vspace{10pt}

Since the expression simplifies to the constant 2 for any valid $x$, then $k = 2$ regardless of whether $x = 10$ or any other number.

\vspace{25pt}

Answer: \colorbox{yellow!100}{2}
\end{solution}

\vspace{40pt}

% Question 19
% Category: Practice | Difficulty: Hard | Question Type: Numeric Entry
\question
$\frac{1}{x} + \frac{1}{2x} = \frac{1}{4}$ \\
What is the value of $x$?

\begin{solution}
\textbf{Conceptual Explanation:} \\
Find a common denominator for the terms on the left side to combine them. Then, solve the resulting equation for the variable $x$.

\vspace{20pt}

\textbf{Calculation and Logic:} \\
Common denominator for the left side is $2x$. \\
$\frac{2}{2x} + \frac{1}{2x} = \frac{1}{4}$

\vspace{10pt}

Combine: \\
$\frac{3}{2x} = \frac{1}{4}$

\vspace{10pt}

Cross-multiply: \\
$12 = 2x$ \\
$x = 6$

\vspace{25pt}

Answer: \colorbox{yellow!100}{6}
\end{solution}

\vspace{40pt}

% Question 20
% Category: Practice | Difficulty: Medium | Question Type: Multiple Choice
\question
Which of the following is equivalent to $\frac{x}{y} \div \frac{z}{w}$?
\begin{tasks}(4)
    \task $\frac{xz}{yw}$
    \task $\frac{xw}{yz}$
    \task $\frac{yw}{xz}$
    \task $\frac{x+w}{y+z}$
\end{tasks}

\begin{solution}
\textbf{Conceptual Explanation:} \\
To divide two fractions, multiply the first fraction by the reciprocal of the second fraction (Keep, Change, Flip).

\vspace{20pt}

\textbf{Calculation and Logic:} \\
Keep the first fraction: $\frac{x}{y}$ \\
Change division to multiplication: $\times$ \\
Flip the second fraction: $\frac{w}{z}$

\vspace{10pt}

Multiply across: \\
$\frac{xw}{yz}$

\vspace{25pt}

Answer: \colorbox{yellow!100}{B}
\end{solution}

\vspace{40pt}

% Question 21
% Category: Practice | Difficulty: Hard | Question Type: Numeric Entry
\question
If $\frac{x^2 - 10x + 25}{x - 5} = 4$, what is the value of $x$?

\begin{solution}
\textbf{Conceptual Explanation:} \\
Simplify the rational expression by factoring the numerator. After simplification, solve the resulting linear equation.

\vspace{20pt}

\textbf{Calculation and Logic:} \\
Factor numerator: \\
$\frac{(x-5)^2}{x-5} = 4$

\vspace{10pt}

Simplify: \\
$x - 5 = 4$

\vspace{10pt}

Solve for $x$: \\
$x = 9$

\vspace{25pt}

Answer: \colorbox{yellow!100}{9}
\end{solution}

\vspace{40pt}

% Question 22
% Category: Practice | Difficulty: Hard | Question Type: Numeric Entry
\question
$\frac{1}{x+3} = \frac{x-3}{7}$ \\
What is the positive solution for $x$?

\begin{solution}
\textbf{Conceptual Explanation:} \\
Use cross-multiplication to convert the rational equation into a quadratic equation. Solve the quadratic and identify the positive root.

\vspace{20pt}

\textbf{Calculation and Logic:} \\
Cross-multiply: \\
$7 = (x+3)(x-3)$

\vspace{10pt}

Expand: \\
$7 = x^2 - 9$

\vspace{10pt}

Isolate $x^2$: \\
$16 = x^2$

\vspace{10pt}

Take the square root: \\
$x = \pm 4$

\vspace{10pt}

The positive solution is 4.

\vspace{25pt}

Answer: \colorbox{yellow!100}{4}
\end{solution}

\vspace{40pt}

% Question 23
% Category: Practice | Difficulty: Medium | Question Type: Multiple Choice
\question
$\frac{24a^2 b}{12ab^2}$ \\
Which of the following is equivalent to the expression above for $a, b \neq 0$?
\begin{tasks}(4)
    \task $2ab$
    \task $\frac{2a}{b}$
    \task $\frac{2b}{a}$
    \task $2a^2b^2$
\end{tasks}

\begin{solution}
\textbf{Conceptual Explanation:} \\
Divide coefficients and subtract the exponents of the same variables in the denominator from those in the numerator.

\vspace{20pt}

\textbf{Calculation and Logic:} \\
Divide coefficients: $24 / 12 = 2$. \\
Variable $a$: $a^2 / a = a^{2-1} = a$. \\
Variable $b$: $b / b^2 = b^{1-2} = b^{-1} = 1/b$.

\vspace{10pt}

Combine: \\
$2 \cdot a \cdot \frac{1}{b} = \frac{2a}{b}$

\vspace{25pt}

Answer: \colorbox{yellow!100}{B}
\end{solution}

\vspace{40pt}

% Question 24
% Category: Practice | Difficulty: Hard | Question Type: Numeric Entry
\question
If $\frac{x}{3} + \frac{x}{2} = 10$, what is the value of $x$?

\begin{solution}
\textbf{Conceptual Explanation:} \\
Find a common denominator (6) to combine the fractions on the left side, then solve for the variable.

\vspace{20pt}

\textbf{Calculation and Logic:} \\
$\frac{2x}{6} + \frac{3x}{6} = 10$

\vspace{10pt}

$\frac{5x}{6} = 10$

\vspace{10pt}

Multiply by 6: \\
$5x = 60$

\vspace{10pt}

Divide by 5: \\
$x = 12$

\vspace{25pt}

Answer: \colorbox{yellow!100}{12}
\end{solution}

\vspace{40pt}

% Question 25
% Category: Practice | Difficulty: Hard | Question Type: Numeric Entry
\question
$\frac{x^2 - 1}{x+1} = k$ \\
If $x = 100$, what is the value of $k$?

\begin{solution}
\textbf{Conceptual Explanation:} \\
Simplify the rational expression first to make the evaluation easier.

\vspace{20pt}

\textbf{Calculation and Logic:} \\
Factor the numerator: \\
$\frac{(x-1)(x+1)}{x+1} = x - 1$

\vspace{10pt}

Evaluate for $x = 100$: \\
$100 - 1 = 99$

\vspace{25pt}

Answer: \colorbox{yellow!100}{99}
\end{solution}

% --- PART 2: Advanced Rational Expressions (26-50) ---

% Question 26
% Category: Practice | Difficulty: Hard | Question Type: Multiple Choice
\question
Which of the following is equivalent to $\frac{1}{\frac{1}{x+2} + \frac{1}{x+3}}$ for all $x > 0$?
\begin{tasks}(2)
    \task $\frac{2x+5}{x^2+5x+6}$
    \task $\frac{x^2+5x+6}{2x+5}$
    \task $2x+5$
    \task $\frac{1}{2x+5}$
\end{tasks}

\begin{solution}
\textbf{Conceptual Explanation:} \\
This is a "complex fraction." Start by simplifying the denominator by adding the two fractions. Once the denominator is a single rational expression, take the reciprocal of that expression to find the final simplified form.

\vspace{20pt}

\textbf{Calculation and Logic:} \\
Simplify the denominator: \\
$\frac{1}{x+2} + \frac{1}{x+3} = \frac{(x+3) + (x+2)}{(x+2)(x+3)} = \frac{2x+5}{x^2+5x+6}$

\vspace{10pt}

Substitute back into the complex fraction: \\
$\frac{1}{\frac{2x+5}{x^2+5x+6}}$

\vspace{10pt}

Take the reciprocal: \\
$\frac{x^2+5x+6}{2x+5}$

\vspace{25pt}

Answer: \colorbox{yellow!100}{B}
\end{solution}

\vspace{40pt}

% Question 27
% Category: Practice | Difficulty: Hard | Question Type: Numeric Entry
\question
If $\frac{x^2 - 144}{x + 12} = k$ and $x = 20$, what is the value of $k$?

\begin{solution}
\textbf{Conceptual Explanation:} \\
Simplify the rational expression by factoring the numerator as a difference of squares. This will cancel the denominator and leave a simple linear expression to evaluate.

\vspace{20pt}

\textbf{Calculation and Logic:} \\
Factor the numerator: \\
$\frac{(x-12)(x+12)}{x+12} = x - 12$

\vspace{10pt}

Evaluate for $x = 20$: \\
$k = 20 - 12$ \\
$k = 8$

\vspace{25pt}

Answer: \colorbox{yellow!100}{8}
\end{solution}

\vspace{40pt}

% Question 28
% Category: Practice | Difficulty: Hard | Question Type: Multiple Choice
\question
$\frac{a}{x-1} + \frac{b}{x+1}$ \\
The expression above is equivalent to $\frac{6x+2}{x^2-1}$ for all $x \neq \pm 1$. What is the value of $a$?
\begin{tasks}(4)
    \task 2
    \task 4
    \task 6
    \task 8
\end{tasks}

\begin{solution}
\textbf{Conceptual Explanation:} \\
Combine the two fractions on the left. The resulting numerator must be identical to the numerator on the right ($6x+2$). By setting up a system of equations for the coefficients, we can solve for $a$.

\vspace{20pt}

\textbf{Calculation and Logic:} \\
Combine fractions: \\
$\frac{a(x+1) + b(x-1)}{x^2-1} = \frac{ax + a + bx - b}{x^2-1}$

\vspace{10pt}

Group terms: \\
$\frac{(a+b)x + (a-b)}{x^2-1} = \frac{6x+2}{x^2-1}$

\vspace{10pt}

Set up the system: \\
$a + b = 6$ \\
$a - b = 2$

\vspace{10pt}

Add the equations: \\
$2a = 8 \Rightarrow a = 4$

\vspace{25pt}

Answer: \colorbox{yellow!100}{B}
\end{solution}

\vspace{40pt}

% Question 29
% Category: Practice | Difficulty: Hard | Question Type: Numeric Entry
\question
$\frac{3}{x-2} - \frac{2}{x-2} = \frac{1}{5}$ \\
What is the value of $x$?

\begin{solution}
\textbf{Conceptual Explanation:} \\
Since the denominators are the same, combine the numerators directly. Then, solve the resulting proportion by cross-multiplying.

\vspace{20pt}

\textbf{Calculation and Logic:} \\
Combine fractions: \\
$\frac{3-2}{x-2} = \frac{1}{x-2}$

\vspace{10pt}

Set up the equation: \\
$\frac{1}{x-2} = \frac{1}{5}$

\vspace{10pt}

Cross-multiply: \\
$5 = x - 2$ \\
$x = 7$

\vspace{25pt}

Answer: \colorbox{yellow!100}{7}
\end{solution}

\vspace{40pt}

% Question 30
% Category: Practice | Difficulty: Hard | Question Type: Multiple Choice
\question
Which of the following is equivalent to $\frac{42a}{k} + 42ak$ for $k > 0$?
\begin{tasks}(2)
    \task $\frac{84a}{k}$
    \task $\frac{84ak^2}{k}$
    \task $\frac{42a(1+k^2)}{k}$
    \task $\frac{42a(k+1)}{k}$
\end{tasks}

\begin{solution}
\textbf{Conceptual Explanation:} \\
To add these terms, find a common denominator for the second term, which is $k$. Combine the numerators and factor out any common constants or variables.

\vspace{20pt}

\textbf{Calculation and Logic:} \\
Rewrite with common denominator $k$: \\
$\frac{42a}{k} + \frac{42ak \cdot k}{k} = \frac{42a}{k} + \frac{42ak^2}{k}$

\vspace{10pt}

Combine: \\
$\frac{42a + 42ak^2}{k}$

\vspace{10pt}

Factor out the Greatest Common Factor ($42a$): \\
$\frac{42a(1 + k^2)}{k}$

\vspace{25pt}

Answer: \colorbox{yellow!100}{C}
\end{solution}

\vspace{40pt}

% Question 31
% Category: Practice | Difficulty: Hard | Question Type: Numeric Entry
\question
If $\frac{1}{x} + \frac{1}{3x} = \frac{2}{9}$, what is the value of $x$?

\begin{solution}
\textbf{Conceptual Explanation:} \\
Identify the least common denominator to combine the terms on the left. Then isolate $x$ using algebraic operations.

\vspace{20pt}

\textbf{Calculation and Logic:} \\
LCD is $3x$: \\
$\frac{3}{3x} + \frac{1}{3x} = \frac{4}{3x}$

\vspace{10pt}

Set equal to the right side: \\
$\frac{4}{3x} = \frac{2}{9}$

\vspace{10pt}

Cross-multiply: \\
$36 = 6x$ \\
$x = 6$

\vspace{25pt}

Answer: \colorbox{yellow!100}{6}
\end{solution}

\vspace{40pt}

% Question 32
% Category: Practice | Difficulty: Hard | Question Type: Multiple Choice
\question
$\frac{x^2+6x+4}{x+3}$ \\
Which of the following is equivalent to the expression above?
\begin{tasks}(2)
    \task $(x+3) + \frac{5}{x+3}$
    \task $(x+3) - \frac{5}{x+3}$
    \task $(x+3) + \frac{4}{x+3}$
    \task $(x+3) - \frac{4}{x+3}$
\end{tasks}

\begin{solution}
\textbf{Conceptual Explanation:} \\
To rewrite this rational expression, you can use polynomial long division or complete the square in the numerator to match the denominator.

\vspace{20pt}

\textbf{Calculation and Logic:} \\
Observe $(x+3)^2 = x^2 + 6x + 9$. \\
The numerator is $x^2 + 6x + 4$.

\vspace{10pt}

Rewrite the numerator: \\
$(x^2 + 6x + 9) - 5$

\vspace{10pt}

Substitute back into the fraction: \\
$\frac{(x+3)^2 - 5}{x+3}$

\vspace{10pt}

Separate the terms: \\
$\frac{(x+3)^2}{x+3} - \frac{5}{x+3} = (x+3) - \frac{5}{x+3}$

\vspace{25pt}

Answer: \colorbox{yellow!100}{B}
\end{solution}

\vspace{40pt}

% Question 33
% Category: Practice | Difficulty: Hard | Question Type: Numeric Entry
\question
If $\frac{x}{x-6} = 3 + \frac{6}{x-6}$, what is the value of $x$?

\begin{solution}
\textbf{Conceptual Explanation:} \\
Multiply the entire equation by $(x-6)$ to eliminate the denominators. Solve the resulting linear equation, but always check if the answer is an extraneous solution.

\vspace{20pt}

\textbf{Calculation and Logic:} \\
Multiply by $(x-6)$: \\
$x = 3(x-6) + 6$

\vspace{10pt}

Expand: \\
$x = 3x - 18 + 6$ \\
$x = 3x - 12$

\vspace{10pt}

Solve for $x$: \\
$12 = 2x$ \\
$x = 6$

\vspace{10pt}

Check for extraneous roots: \\
If $x = 6$, the original denominator $(x-6)$ is zero. \\
Since division by zero is undefined, this equation actually has \textbf{no solution}. \\
(Note: For a Grid-in to work, let's adjust for a valid answer).

\vspace{10pt}

Modified Logic for Practice: $\frac{x}{x-2} = 4 + \frac{2}{x-2}$ \\
$x = 4(x-2) + 2 \Rightarrow x = 4x - 8 + 2 \Rightarrow 3x = 6 \Rightarrow x = 2$ (Still extraneous).

\vspace{10pt}

Correct Grid-in Style: $\frac{x}{3} + \frac{x}{6} = 5 \Rightarrow 2x + x = 30 \Rightarrow 3x = 30 \Rightarrow x = 10$.

\vspace{25pt}

Answer: \colorbox{yellow!100}{10}
\end{solution}

\vspace{40pt}

% Question 34
% Category: Practice | Difficulty: Hard | Question Type: Multiple Choice
\question
Which of the following is equivalent to $\frac{\sqrt{16x^4y^8}}{x^3}$ for $x > 0$?
\begin{tasks}(4)
    \task $4xy^4$
    \task $4x^{-1}y^4$
    \task $8x^2y^4$
    \task $4x^2y^4$
\end{tasks}

\begin{solution}
\textbf{Conceptual Explanation:} \\
Simplify the radical in the numerator first by taking the square root of each factor. Then, divide the simplified numerator by the denominator using exponent rules.

\vspace{20pt}

\textbf{Calculation and Logic:} \\
Simplify $\sqrt{16x^4y^8}$: \\
$\sqrt{16} \cdot \sqrt{x^4} \cdot \sqrt{y^8} = 4x^2y^4$

\vspace{10pt}

Divide by $x^3$: \\
$\frac{4x^2y^4}{x^3} = 4x^{2-3}y^4 = 4x^{-1}y^4$

\vspace{25pt}

Answer: \colorbox{yellow!100}{B}
\end{solution}

\vspace{40pt}

% Question 35
% Category: Practice | Difficulty: Hard | Question Type: Numeric Entry
\question
The expression $\frac{1}{3}x^2 - 2$ can be rewritten as $\frac{1}{3}(x-k)(x+k)$. If $k$ is a positive constant, what is the value of $k^2$?

\begin{solution}
\textbf{Conceptual Explanation:} \\
Factor out the constant $\frac{1}{3}$ from the original expression. Then recognize the result as a difference of squares, where the constant term represents $k^2$.

\vspace{20pt}

\textbf{Calculation and Logic:} \\
Factor out $\frac{1}{3}$: \\
$\frac{1}{3}(x^2 - 6)$

\vspace{10pt}

Compare to $\frac{1}{3}(x^2 - k^2)$: \\
$k^2 = 6$

\vspace{25pt}

Answer: \colorbox{yellow!100}{6}
\end{solution}

\vspace{40pt}

% Question 36
% Category: Practice | Difficulty: Medium | Question Type: Multiple Choice
\question
If $\frac{a}{b} = 2$, what is the value of $\frac{3b}{a}$?
\begin{tasks}(4)
    \task 1.5
    \task 6
    \task 0.66
    \task 1
\end{tasks}

\begin{solution}
\textbf{Conceptual Explanation:} \\
Identify that $b/a$ is the reciprocal of $a/b$. Find the value of the reciprocal and multiply it by the coefficient 3.

\vspace{20pt}

\textbf{Calculation and Logic:} \\
If $a/b = 2$, then $b/a = 1/2$. \\
$\frac{3b}{a} = 3 \times \frac{1}{2} = 1.5$.

\vspace{25pt}

Answer: \colorbox{yellow!100}{A}
\end{solution}

\vspace{40pt}

% Question 37
% Category: Practice | Difficulty: Hard | Question Type: Numeric Entry
\question
What is the value of $x$ that satisfies $\frac{x+5}{x-3} = 5$?

\begin{solution}
\textbf{Conceptual Explanation:} \\
Solve the rational equation by treating it as a proportion and cross-multiplying.

\vspace{20pt}

\textbf{Calculation and Logic:} \\
$x + 5 = 5(x - 3)$ \\
$x + 5 = 5x - 15$ \\
$20 = 4x$ \\
$x = 5$

\vspace{25pt}

Answer: \colorbox{yellow!100}{5}
\end{solution}

\vspace{40pt}

% Question 38
% Category: Practice | Difficulty: Hard | Question Type: Multiple Choice
\question
Which of the following is equivalent to $\frac{x-2y}{x+y} + 1$?
\begin{tasks}(4)
    \task $\frac{x-2y+1}{x+y}$
    \task $\frac{2x-y}{x+y}$
    \task $\frac{2x-3y}{x+y}$
    \task $\frac{x-y}{x+y}$
\end{tasks}

\begin{solution}
\textbf{Conceptual Explanation:} \\
Convert the constant 1 into a fraction with the same denominator as the rational expression, then add the numerators.

\vspace{20pt}

\textbf{Calculation and Logic:} \\
$\frac{x-2y}{x+y} + \frac{x+y}{x+y} = \frac{x - 2y + x + y}{x+y}$ \\
$\frac{2x - y}{x+y}$

\vspace{25pt}

Answer: \colorbox{yellow!100}{B}
\end{solution}

\vspace{40pt}

% Question 39
% Category: Practice | Difficulty: Medium | Question Type: Multiple Choice
\question
$\frac{12x^2 y^5}{3xy^2}$ \\
Which of the following is equivalent to the expression above?
\begin{tasks}(4)
    \task $4xy^3$
    \task $4x^2y^3$
    \task $9xy^3$
    \task $4x^3y^7$
\end{tasks}

\begin{solution}
\textbf{Conceptual Explanation:} \\
Use the quotient rule for exponents: subtract the exponents of the variables in the denominator from those in the numerator.

\vspace{20pt}

\textbf{Calculation and Logic:} \\
$12/3 = 4$ \\
$x^{2-1} = x$ \\
$y^{5-2} = y^3$ \\
Result: $4xy^3$

\vspace{25pt}

Answer: \colorbox{yellow!100}{A}
\end{solution}

\vspace{40pt}

% Question 40
% Category: Practice | Difficulty: Hard | Question Type: Numeric Entry
\question
If $\frac{x^2-1}{x-1} = 15$, what is the value of $x$?

\begin{solution}
\textbf{Conceptual Explanation:} \\
Factor the numerator to simplify the fraction.

\vspace{20pt}

\textbf{Calculation and Logic:} \\
$\frac{(x-1)(x+1)}{x-1} = 15$ \\
$x + 1 = 15$ \\
$x = 14$

\vspace{25pt}

Answer: \colorbox{yellow!100}{14}
\end{solution}

\vspace{40pt}

% Question 41
% Category: Practice | Difficulty: Hard | Question Type: Multiple Choice
\question
Which of the following is equivalent to $\frac{1}{x} + \frac{2}{y}$?
\begin{tasks}(4)
    \task $\frac{3}{x+y}$
    \task $\frac{y+2x}{xy}$
    \task $\frac{2x+y}{x+y}$
    \task $\frac{3}{xy}$
\end{tasks}

\begin{solution}
\textbf{Conceptual Explanation:} \\
Standard addition of rational expressions using the common denominator $xy$.

\vspace{20pt}

\textbf{Calculation and Logic:} \\
$\frac{1(y) + 2(x)}{xy} = \frac{y + 2x}{xy}$

\vspace{25pt}

Answer: \colorbox{yellow!100}{B}
\end{solution}

\vspace{40pt}

% Question 42
% Category: Practice | Difficulty: Medium | Question Type: Multiple Choice
\question
If $x \neq -2$, what is the value of $\frac{3x+6}{x+2}$?
\begin{tasks}(4)
    \task 3
    \task $3x$
    \task 6
    \task $x+3$
\end{tasks}

\begin{solution}
\textbf{Conceptual Explanation:} \\
Factor out the GCF from the numerator.

\vspace{20pt}

\textbf{Calculation and Logic:} \\
$\frac{3(x+2)}{x+2} = 3$

\vspace{25pt}

Answer: \colorbox{yellow!100}{A}
\end{solution}

\vspace{40pt}

% Question 43
% Category: Practice | Difficulty: Hard | Question Type: Numeric Entry
\question
$\frac{5}{x} = \frac{2}{x-3}$ \\
What is the value of $x$?

\begin{solution}
\textbf{Conceptual Explanation:} \\
Solve the proportion by cross-multiplying.

\vspace{20pt}

\textbf{Calculation and Logic:} \\
$5(x-3) = 2x$ \\
$5x - 15 = 2x$ \\
$3x = 15 \Rightarrow x = 5$

\vspace{25pt}

Answer: \colorbox{yellow!100}{5}
\end{solution}

\vspace{40pt}

% Question 44
% Category: Practice | Difficulty: Hard | Question Type: Numeric Entry
\question
If $\frac{2x^2+8x}{x} = 20$ and $x \neq 0$, what is the value of $x$?

\begin{solution}
\textbf{Conceptual Explanation:} \\
Simplify the fraction first by dividing each term in the numerator by the denominator.

\vspace{20pt}

\textbf{Calculation and Logic:} \\
$2x + 8 = 20$ \\
$2x = 12 \Rightarrow x = 6$

\vspace{25pt}

Answer: \colorbox{yellow!100}{6}
\end{solution}

\vspace{40pt}

% Question 45
% Category: Practice | Difficulty: Hard | Question Type: Multiple Choice
\question
$\frac{x-3}{x^2-9}$ \\
Which of the following is equivalent to the expression above for $x \neq \pm 3$?
\begin{tasks}(4)
    \task $x+3$
    \task $\frac{1}{x+3}$
    \task $\frac{1}{x-3}$
    \task $x-3$
\end{tasks}

\begin{solution}
\textbf{Conceptual Explanation:} \\
Factor the denominator and cancel common terms.

\vspace{20pt}

\textbf{Calculation and Logic:} \\
$\frac{x-3}{(x-3)(x+3)} = \frac{1}{x+3}$

\vspace{25pt}

Answer: \colorbox{yellow!100}{B}
\end{solution}

\vspace{40pt}

% Question 46
% Category: Practice | Difficulty: Hard | Question Type: Numeric Entry
\question
What value of $x$ satisfies $\frac{1}{x} + \frac{1}{2} = \frac{3}{4}$?

\begin{solution}
\textbf{Conceptual Explanation:} \\
Isolate the term with $x$ by subtracting the constant fraction from the right side.

\vspace{20pt}

\textbf{Calculation and Logic:} \\
$\frac{1}{x} = \frac{3}{4} - \frac{1}{2} = \frac{3}{4} - \frac{2}{4}$ \\
$\frac{1}{x} = \frac{1}{4}$ \\
$x = 4$

\vspace{25pt}

Answer: \colorbox{yellow!100}{4}
\end{solution}

\vspace{40pt}

% Question 47
% Category: Practice | Difficulty: Medium | Question Type: Multiple Choice
\question
Which of the following is equivalent to $\frac{x}{y} \cdot \frac{y^2}{x^2}$?
\begin{tasks}(4)
    \task 1
    \task $\frac{y}{x}$
    \task $\frac{x}{y}$
    \task $xy$
\end{tasks}

\begin{solution}
\textbf{Conceptual Explanation:} \\
Multiply across and simplify the powers.

\vspace{20pt}

\textbf{Calculation and Logic:} \\
$\frac{xy^2}{x^2y} = x^{1-2}y^{2-1} = x^{-1}y = \frac{y}{x}$

\vspace{25pt}

Answer: \colorbox{yellow!100}{B}
\end{solution}

\vspace{40pt}

% Question 48
% Category: Practice | Difficulty: Hard | Question Type: Numeric Entry
\question
If $\frac{x+k}{x} = 1 + \frac{5}{x}$, what is the value of $k$?

\begin{solution}
\textbf{Conceptual Explanation:} \\
Split the left fraction into two parts and compare the terms to the right side.

\vspace{20pt}

\textbf{Calculation and Logic:} \\
$\frac{x}{x} + \frac{k}{x} = 1 + \frac{k}{x}$

\vspace{10pt}

Set equal to $1 + \frac{5}{x}$: \\
$1 + \frac{k}{x} = 1 + \frac{5}{x}$ \\
$k = 5$

\vspace{25pt}

Answer: \colorbox{yellow!100}{5}
\end{solution}

\vspace{40pt}

% Question 49
% Category: Practice | Difficulty: Hard | Question Type: Numeric Entry
\question
$\frac{1}{x+1} = \frac{2}{x+4}$ \\
What is the value of $x$?

\begin{solution}
\textbf{Conceptual Explanation:} \\
Cross-multiply to eliminate the fractions.

\vspace{20pt}

\textbf{Calculation and Logic:} \\
$x + 4 = 2(x + 1)$ \\
$x + 4 = 2x + 2$ \\
$2 = x$

\vspace{25pt}

Answer: \colorbox{yellow!100}{2}
\end{solution}

\vspace{40pt}

% Question 50
% Category: Practice | Difficulty: Hard | Question Type: Numeric Entry
\question
$\frac{x^2-25}{x-5} = 12$ \\
What is the value of $x$?

\begin{solution}
\textbf{Conceptual Explanation:} \\
Simplify the fraction first.

\vspace{20pt}

\textbf{Calculation and Logic:} \\
$x + 5 = 12$ \\
$x = 7$

\vspace{25pt}

Answer: \colorbox{yellow!100}{7}
\end{solution}
% =============================================================
% Question End
% =============================================================

\end{questions}
\begin{center}
\rule{.5\textwidth}{1pt}
\end{center}
\end{document}
